\needspace{10\baselineskip}
\begin{center}
\pdfbookmark[1]{7. On Art. 90 and 100.}{zusatz-007}
[7. On Art. 90 and 100.]
\end{center}
\begin{center}
\begin{tabular}{c}\hline
\small Astronomisches Jahrbuch für 1814, p. 256. Berlin 1811.\\\hline
\end{tabular}
\end{center}
.... In accordance with your wish, I also append a small addition to my Theoria motus corporum coelestium.

To solve the important problem of determining the elliptic or hyperbolic elements from two radius vectors and the included angle, I have made very advantageous use of an auxiliary quantity, \(\xi\) for the ellipse and \(\zeta\) for the hyperbola, for which I appended a table to that work. This table is calculated from a continued fraction given there, whose complete derivation, however, is not given there, and whose theoretical development, which is connected with other investigations, I have not yet had occasion to present. Many will therefore be glad to find another method indicated here, by which that auxiliary quantity could have been calculated just as conveniently.

We have (Art. 90)
\[
\xi=x-\tfrac56+\frac{10}{9X}=\frac{xX-\tfrac56X+\tfrac{10}{9}}{X}\mathpunct{\cdot}
\]

\SourceLeaf{307}
The numerator of this fraction is readily transformed, on substituting for \(X\) the series given there, into
\[
\begin{aligned}
&\tfrac8{105}xx\Bigl(1+\frac{2.8}{9}x+\frac{3.8.10}{9.11}xx+\frac{4.8.10.12}{9.11.13}x^3\\
&\qquad+\frac{5.8.10.12.14}{9.11.13.15}x^4+\text{etc.}\Bigr).
\end{aligned}
\]

If, therefore, we set the series
\[
1+\frac{2.8}{9}x+\frac{3.8.10}{9.11}xx+\ldots=A,
\]
we obtain
\[
\begin{aligned}
xX-\tfrac56X+\tfrac{10}{9}&=\tfrac8{105}Axx\\[.5ex]
X&=\frac{\tfrac43\bigl(1-\tfrac{12}{175}Axx\bigr)}{1-\tfrac65x}\\[.5ex]
\xi&=\frac{\tfrac2{35}Axx\bigl(1-\tfrac65x\bigr)}{1-\tfrac{12}{175}Axx}\mathpunct{\raisebox{.45ex}{,}}
\end{aligned}
\]
by which formula \(\xi\) can always be calculated conveniently and reliably. For \(\zeta\), one need only put \(z\) in place of \(x\).

I only add that \(A\) can be calculated still more conveniently by the following formula
\[
\begin{aligned}
A=(1-x)^{-\frac32}\Bigl(&1+\frac{1.5}{2.9}x+\frac{1.3.5.7}{2.4.9.11}xx\\
&+\frac{1.3.5.5.7.9}{2.4.6.9.11.13}x^3+\text{etc.}\Bigr);
\end{aligned}
\]
but the derivation of this series from the preceding one rests on grounds that cannot be developed here\HistoricalNoteEight{Cf. Volume III p. 209.]}.
\begin{center}\rule{.18\textwidth}{.35pt}\end{center}
